\documentclass{beamer}
\usetheme{zgca}

% Set the background image for the title slide (recolored building image)
\titlebackground{assets/background}

\title{中关村学院 Beamer 模板展示}
\subtitle{使用 LaTeX 制作学术报告演示文稿}
\course{中关村学院卓越学术报告系列}
\author{报告人：张三 \quad 导师：李四}
\IDnumber{ZGCA-2026001}
\date{2026 年 6 月 12 日}

\begin{document}
\maketitle

\begin{frame}{模板简介}
\framesubtitle{基于 ZGCA 视觉识别系统的学术幻灯片}

本模板是专门为中关村学院学术报告定制的 Beamer 演示文稿模板，适配了学院的视觉标识（主题色：\texttt{\#0363c9}）。

\vspace{\baselineskip}

主要适配内容包括：
\begin{itemize}
    \item 结合中关村学院与中关村人工智能研究院的联合 logo 标识
    \item 主题色变更为标准的中关村学院蓝色（\texttt{\#0363c9}）
    \item 封面图和过渡页提取自学院 C5 大楼，并保留了原本的学院色彩与排版样式
    \item 内页页眉及页脚设计，包含学院纸飞机 icon 及双 logo 边角装饰
\end{itemize}

\end{frame}

\section{基本元素演示}

\begin{frame}[fragile]{列表与强调}
\framesubtitle{演示如何展示列表和高亮内容}

我们可以方便地使用多种格式的列表和强调：
\begin{itemize}
    \item 这是一个列表项，支持在演示中逐步淡入展现
    \item 推荐使用 \verb|\alert{内容}| 来高亮需要引人注意的词语，例如：\alert{中关村学院}
    \item 主题色相关的颜色宏定义在 \texttt{zgcacolor.sty} 包中，主要包括：
    \begin{itemize}
        \item 主题色：\testcolor{maincolor}
        \item 辅助黄色：\testcolor{zgcayellow}（适用于 alert 文字）
        \item 浅底填充色：\testcolor{zgcagrey}（适用于块环境背景）
    \end{itemize}
\end{itemize}
\end{frame}

\begin{frame}{学术数学公式支持}
\framesubtitle{完美契合 LaTeX 的数学公式排版}

Beamer 模板对数学公式的支持极其强大，模板自动集成了 serif 字体用于公式排版以确保其美观易读：

\begin{equation}
\mathrm{i}\,\hslash\frac{\partial}{\partial t} \Psi(\mathbf{r},t) =
-\frac{\hslash^2}{2\,m}\nabla^2\Psi(\mathbf{r},t)
+ V(\mathbf{r})\Psi(\mathbf{r},t)
\end{equation}

\begin{block}{深度学习中的记忆与持续学习机制}
持续学习（Continual Learning）旨在克服人工神经网络在学习新任务时对旧任务的“灾难性遗忘”（Catastrophic Forgetting）。数学表述可由正则化项引入：
\begin{equation*}
\mathcal{L}(\theta) = \mathcal{L}_B(\theta) + \sum_{i} \frac{\lambda}{2} F_i (\theta_i - \theta_{A,i}^*)^2
\end{equation*}
\end{block}

\end{frame}

\begin{frame}[fragile]{卡片块设计 (Blocks)}
\framesubtitle{利用块结构来区分不同类型的内容}

\begin{columns}
\begin{column}{0.48\textwidth}
\begin{block}{标准块 (Block)}
这是最常用的块结构。块的头部使用主题色（标准蓝色），内容背景使用微蓝色的淡灰色。
\small
\begin{verbatim}
\begin{block}{标题}
内容...
\end{block}
\end{verbatim}
\end{block}
\end{column}

\begin{column}{0.48\textwidth}
\begin{colorblock}[white]{maincolor}{全色块 (Colorblock)}
头部和内容使用统一的主题色背景，文字为白色。适合放置一些非常核心的结论或公式。
\small
\begin{verbatim}
\begin{colorblock}{颜色}{标题}
内容...
\end{colorblock}
\end{verbatim}
\end{colorblock}
\end{column}
\end{columns}
\end{frame}

\section{高级版式与个人化}

\footlinecolor{maincolor}
\begin{frame}[fragile]{切换底部栏底色}
\framesubtitle{为特定页面改变版式风格}

\begin{itemize}
    \item 默认情况下，幻灯片采用透明底部栏，页面清爽大气
    \item 可以使用 \verb|\footlinecolor{颜色}| 来为特定的页面添加带底色的底部栏
    \item 例如本页的底部栏已经被设置为了：\testcolor{maincolor}
    \item 若要恢复默认的无底色状态，请在下一页的 \verb|frame| 之前执行：\verb|\footlinecolor{}|
\end{itemize}

\begin{block}{注意事项}
请勿在 \verb|themecolor{main}|（即全黑/全主题色幻灯片）中改变页脚颜色。
\end{block}
\end{frame}

\footlinecolor{}
\begin{chapter}[assets/background]{maincolor}{过渡章节页展示}
\begin{itemize}
    \item 过渡章节页采用全主题色背景设计
    \item 右侧可以使用半截裁剪的背景图进行点缀
    \item 这一页采用 \texttt{\textbackslash begin\string{chapter\string}[图片]\string{背景色\string}\string{标题\string}} 环境创建
\end{itemize}
\end{chapter}

\begin{sidepic}{assets/background}{侧图版式演示}
\begin{itemize}
    \item 侧图版式提供了左侧文字、右侧整版插图的经典设计
    \item 这一页采用 \texttt{\textbackslash begin\string{sidepic\string}\string{图片\string}\string{标题\string}} 环境创建
    \item 适合在幻灯片中插入相关的实物、图表或风景图
\end{itemize}
\end{sidepic}

\begin{frame}{总结与展望}
\framesubtitle{演示如何使用表格排版总结数据}

\begin{table}
\centering
\begin{tabular}{lcc}
\hline
\textbf{指标} & \textbf{传统 PPT} & \textbf{ZGCA Beamer (本模板)} \\
\hline
数学公式美观度 & 一般 & \textbf{极佳} \\
视觉色彩一致性 & 难保持 & \textbf{优秀 (适配 \#0363c9)} \\
排版自动控制 & 手动调整 & \textbf{自动排版} \\
学院视觉契合度 & 低 & \textbf{高度适配} \\
\hline
\end{tabular}
\caption{演示文稿排版工具对比}
\end{table}

\end{frame}

\backmatter
\end{document}
